Para cada entero positivo $n$, definimos la función phi de Euler, denotada por $\Phi(n)$, como el número de enteros positivos menores o iguales que $n$ y primos relativos con $n$. Mediante el principio de inclusión-exclusión demuestre que si $n = \prod_{i=1}^{k} p_i^{\,n_i}$ es la representación canónica de $n$, entonces
\[\Phi(n)=\prod_{i=1}^{k}\bigl(p_i^{\,n_i}-p_i^{\,n_i-1}\bigr)=n\prod_{i=1}^{k}\Bigl(1-\tfrac1{p_i}\Bigr).\]
